% ------------------------------------------------------------------------------
% Estrategia y resultados de validación
% ------------------------------------------------------------------------------

Este capítulo describe la estrategia de pruebas aplicada a \textbf{UniHub}. Se distingue entre pruebas deterministas, que se ejecutan sin depender de fuentes externas, y evaluaciones empíricas, que contrastan el resultado con publicaciones oficiales. Esta separación evita confundir una prueba de software con una medición de cobertura nacional.

\section{Estrategia de Pruebas}
La validación se organiza en cuatro niveles: pruebas unitarias de normalización y reglas de negocio; pruebas de integración entre crawler, persistencia, API y cliente web; pruebas de regresión del parser BOE; y evaluaciones empíricas contra documentos oficiales. Las pruebas con red se ejecutan explícitamente, de forma secuencial y respetando los retardos y las restricciones de acceso de las fuentes consultadas.

\section{Pruebas Unitarias y de Integración}
La suite automatizada cubre de forma exhaustiva los contratos funcionales del sistema mediante \textbf{38 suites de prueba unitarias e integradas}:
\begin{itemize}
    \item \textbf{Ingesta RUCT, BOE y Persistencia Dual:} Clasificación de titulaciones, filtro de vigencia, persistencia transaccional SQLite WAL y recuperación ante corrupción (\nolinkurl{test_checkpoint_resilience.py}, \nolinkurl{test_crawl_ledger_resilience.py}, \nolinkurl{test_cache_recovery.py}, \nolinkurl{test_boe_search_discovery.py}).
    \item \textbf{Políticas de Red, HTTP/2 y Cortesía:} Normalización canónica de URLs, verificación de cabeceras PDF, política \texttt{robots.txt} con seguimiento de redirecciones HTTPS, reintentos exponenciales ante HTTP 429, aislamiento de dominio mediante \emph{Circuit Breaker} y conmutación resiliente ante tramas HTTP/2 anómalas (\nolinkurl{test_downloader_domain_boundaries.py}, \nolinkurl{test_crawler_politeness_concurrency.py}, \nolinkurl{test_robots_policy_redirects.py}, \nolinkurl{test_http_cache_ttl.py}, \nolinkurl{test_downloader_illegal_status_line_recovery.py}).
    \item \textbf{Guías Docentes y Temarios EEES:} Análisis de estructuras tabulares, extracción semántica multilingüe, normalización porcentual del desglose de evaluación, poda inteligente de subdominios ciegos anti-404 y resolución negativa de candidatas en $0\text{ ms}$ (\nolinkurl{test_subject_guide_crawler.py}, \nolinkurl{test_subject_guide_discovery_cache.py}, \nolinkurl{test_guide_pruning_anti_404.py}, \nolinkurl{test_subject_guide_template_learning.py}).
    \item \textbf{Auditoría de Procesos y Concurrencia:} Gestión segura de descriptores nativos C++, recolección de cerrojos huérfanos por marca temporal, cola acotada de telemetría y algoritmos de balanceo de llaves en scripts SPA (\nolinkurl{test_phase1_audit_fixes.py}, \nolinkurl{test_phase1_audit_v2_fixes.py}, \nolinkurl{test_spa_render_optimizations.py}, \nolinkurl{test_spa_interactive_tabs.py}).
    \item \textbf{Cálculo Tarifario y Servicios REST:} Consolidación de precios autonómicos SIIU, matriz ampliada de 51 universidades privadas, carga masiva transaccional ETL, contratos de calidad de planes y endpoints de la API (\nolinkurl{test_plan_publication_api.py}, \nolinkurl{test_quality_persistence_contract.py}, \nolinkurl{test_private_pricing_catalog_expansion.py}).
    \item \textbf{Tablas HTML Complejas y Reconstrucción Matricial:} Propagación de celdas fusionadas por \texttt{rowspan}/\texttt{colspan} e inferencia secuencial de curso lectivo hasta un máximo de 60 ECTS anuales (\nolinkurl{test_course_rowspan_and_ects_inference.py}, \nolinkurl{test_html_internship_rows.py}).
    \item \textbf{Concordancia Curricular y Resumen BOE:} Verificación de candidatos web contra el cuadro oficial de distribución de créditos del BOE (\nolinkurl{test_boe_summary_web_matching.py}, \nolinkurl{test_curriculum_recovery.py}, \nolinkurl{test_course_partitioned_merge.py}).
    \item \textbf{Programas Oficiales de Doctorado (RD 99/2011):} Verificación de los tres patrones universales (listas, subtítulos, tablas), descarte de ruido administrativo, detección multilingüe de escuela y validación curricular sin créditos ECTS lectivos (\nolinkurl{test_doctorate_validation_rd99.py}).
    \item \textbf{Limpieza de Datos y Preparación para Arranque en Frío:} Invariabilidad criptográfica SHA-256 de las semillas maestras, simulación dry-run, purga de artefactos volátiles y arranque limpio en frío sin estados residuales (\nolinkurl{test_cold_start_and_cleanup_readiness.py}).
    \item \textbf{Frontend y Despliegue:} Representación de datos en la SPA, control de accesos RBAC y contratos de orquestación contenerizada (\nolinkurl{test_frontend_data_representation.py}).
\end{itemize}

La ejecución automatizada de la suite completa se realiza mediante el comando:
\begin{lstlisting}[language=bash]
python -m unittest discover -s Codigo/Pruebas -p "test_*.py"
\end{lstlisting}
El resultado de la verificación arroja un \textbf{100\% de pruebas superadas} en las 38 suites sin fallos ni regresiones funcionales, garantizando la solidez matemática y operativa del sistema.

\section{Validación del Parser de Resoluciones BOE}
El parser se valida tanto con fixtures sin red como con publicaciones oficiales. La prueba \texttt{test\_boe\_deep\_parser.py} cubre cuatro contratos críticos de alto riesgo:
\begin{enumerate}
    \item tabla curricular convencional con cabecera explícita;
    \item tabla sin bordes recuperada desde líneas y coordenadas espaciales del PDF;
    \item resolución con dos titulaciones en la misma página, verificando que sólo se conserve la sección solicitada;
    \item aislamiento de la caché en memoria de modelos de documento SHA-256 (\texttt{\_DOCUMENT\_MODEL\_CACHE}) para garantizar independencia entre ejecuciones de prueba.
\end{enumerate}

Además, las pruebas de regresión verifican encabezados multilingües, filas optativas, temporalidad, créditos declarados y documentos multi-titulación. Estas pruebas detectan cambios de comportamiento, pero no miden por sí solas la precisión frente a la fuente publicada.

\subsection{Evaluación Empírica Reproducible}
El script \texttt{boe\_empirical\_evaluation.py} contrasta el PDF consumido por el crawler con la tabla HTML estructurada del BOE. La muestra contiene cinco resoluciones: \texttt{BOE-A-2022-12382}, \texttt{BOE-A-2024-1769}, \texttt{BOE-A-2023-6963}, \texttt{BOE-A-2023-11558} y \texttt{BOE-A-2025-21807}.

En la ejecución de referencia se compararon 181 asignaturas: la precisión fue del 100\%, la cobertura del 100\% y los créditos ECTS declarados coincidieron en los cinco casos. La muestra incluye Grados, Másteres y una resolución multi-titulación. Este resultado evidencia que el parser funciona en los casos ensayados; no debe extrapolarse como una garantía estadística para todos los PDFs históricos o para documentos que no publiquen el detalle curricular.

\section{Pruebas de Rendimiento y Benchmarking}
Para evaluar el impacto de las optimizaciones arquitectónicas en el análisis de documentos PDF, se ejecutó un benchmark comparativo de rendimiento (\texttt{benchmark\_pdf\_optimizations.py}):
\begin{itemize}
    \item \textbf{Pase en Frío (Extracción Geométrica Condicional):} El análisis de un PDF de varias páginas con filtrado previo de páginas irrelevantes toma $\approx 347.2\text{ ms}$.
    \item \textbf{Pase en Caliente (Caché SHA-256 de Modelos de Documento):} Cuando una segunda titulación comparte el mismo documento oficial del BOE, la recuperación instantánea de la estructura geométrica desde memoria RAM toma únicamente $\mathbf{0.92\text{ ms}}$.
    \item \textbf{Ganancia de Eficiencia:} Se obtiene un factor de aceleración de $\mathbf{376.3\times}$ (reducción del 99.7\% del tiempo de CPU en documentos compartidos), permitiendo que la suite completa de 407 tests se ejecute con máxima rapidez y solidez.
\end{itemize}

\section{Pruebas de Ejecución Controlada}
Las ejecuciones limitadas del crawler se emplean para comprobar que el pipeline:
\begin{itemize}
    \item reanuda desde \texttt{checkpoint.json} sin destruir información previa;
    \item registra respuestas, errores y caché en el ledger SQLite;
    \item respeta la política de acceso por dominio y no trata una denegación de \texttt{robots.txt} como error recuperable mediante fuerza bruta;
    \item descarta URLs candidatas inválidas (404/410) en $0\text{ ms}$ a través del índice negativo SQLite WAL;
    \item conserva los resultados incompletos con su estado de fuente, en lugar de rellenarlos con datos no publicados.
\end{itemize}

La duración y la cobertura de una ejecución nacional dependen de las respuestas de RUCT, BOE y las universidades. Por ello, cualquier resultado operativo debe acompañarse del manifiesto de ejecución, el intervalo temporal, los límites de universidades/titulaciones y la configuración de trabajadores utilizada.

\section{Pruebas No Funcionales}
\begin{itemize}
    \item \textbf{Robustez:} cada titulación se procesa de forma aislada; una incidencia de red o de parseo queda registrada sin invalidar el resto de la ejecución.
    \item \textbf{Gestión de Recursos y Fugas:} verificación del cierre explícito de descriptores de sockets, conexiones de bases de datos y punteros nativos C++ en librerías de renderizado.
    \item \textbf{Cortesía:} el crawler consulta \texttt{robots.txt}, aplica retardos por dominio y limita la concurrencia sobre una misma institución.
    \item \textbf{Trazabilidad:} los planes conservan URL, procedencia, fecha de comprobación y, cuando corresponde, la huella SHA-256 del PDF.
    \item \textbf{Rendimiento:} las métricas de tiempo, memoria y caché se registran durante la ejecución. Deben presentarse como mediciones asociadas a un manifiesto, no como constantes universales del sistema.
\end{itemize}

\section{Criterio de Aceptación}
Una versión candidata se acepta cuando la suite determinista finaliza correctamente (100\% de tests superados), la evaluación empírica BOE no reduce precisión ni cobertura respecto a la referencia y las pruebas de integración no detectan roturas de contrato entre crawler, API y frontend. Las titulaciones sin detalle oficial permanecen identificadas como incompletas; la ausencia de datos no se transforma en un dato estimado.
